\BourbakiExerciseGroup

\begin{enumerate}
\def\labelenumi{\arabic{enumi})}
\tightlist
\item
  \begin{enumerate}
  \def\labelenumii{\alph{enumii})}
  \tightlist
  \item
    拓扑群 G 的一个子群 H 按外律 \((s, x) \rightarrow sxs^{-1}\) 逆紧地作用于 G，当且仅当 G 是豪斯多夫的且 H 是紧的（利用命题 2 与命题 4）。
  \end{enumerate}
\end{enumerate}

\begin{enumerate}
\def\labelenumi{\alph{enumi})}
\setcounter{enumi}{1}
\tightlist
\item
  设 G 是一个拓扑群，H 是 G 的一个子群。则映射 \((x,y)\to xy\) （\(H\times G\) 到 G 中的）是逆紧的，当且仅当 H 是拟紧的（考虑 e 在这个映射下的逆像）。
\end{enumerate}

\begin{enumerate}
\def\labelenumi{\arabic{enumi})}
\setcounter{enumi}{1}
\tightlist
\item
  给出一个紧群连续地（因而非逆紧地，参看命题 3）作用于一个非豪斯多夫空间（参看 § 2，习题 13）的例子。
\end{enumerate}

¶ 3) 拓扑群 G 逆紧地作用于齐性空间 G/K，当且仅当 K 是拟紧的。（为证明条件是充分的，试证 G 到 G/K 上的典范映射是逆紧的，并用第一章，§ 10，第 1 节，命题 5。）

\begin{enumerate}
\def\labelenumi{\arabic{enumi})}
\setcounter{enumi}{3}
\item
  给出一个豪斯多夫拓扑群 G 逆紧地作用于一个豪斯多夫拓扑空间 X 的例子，使得映射 \((s, x) \rightarrow s.x\) 不是闭的，且典范映射 \(X \rightarrow X/G\) 不是闭的{[}参看习题 1 b){]}。
\item
  拓扑群 G 的一个子群 H 经左平移逆紧地作用于 G，当且仅当 H 在 G 中是闭的。
\end{enumerate}

* 6) 对每个实数 \(a \geqslant 1\) ，令

\[
\begin{array}{l l} f _ {a} (t) = \left(t + \frac {a (a + 2)}{a + 1}, - \frac {a}{a + 1}\right) & \text {if} \quad t <   - \frac {a}{a + 1}, \\ f _ {a} (t) = (a, t) & \text {if} \quad - \frac {a}{a + 1} \leqslant t \leqslant \frac {a}{a + 1}, \\ f _ {a} (t) = \left(- t + \frac {a (a + 2)}{a + 1}, \frac {a}{a + 1}\right) & \text {if} \quad t > \frac {a}{a + 1}. \end{array}
\]

用 \(C_{a}\) 表示当 t 取遍 R 时所有点 \(f_{a}(t)\) 组成的集合，并设 \(X \subset R^{2}\) 是诸集合 \(C_{a}\) （对 \(a \geqslant 1\) ）与直线 \(D'\) 和 \(D''\) 的并，其中 \(D'\) （相应地，\(D''\) ）是所有点 \((t, -1)\) （相应地，\((t, 1)\) ）（\(t \in R\) ）组成的集合。

\begin{enumerate}
\def\labelenumi{\alph{enumi})}
\tightlist
\item
  群 \(\mathbf{R}\) 作用于 \(\mathbf{X}\) ，按律 \((s, z) \to s.z\) 使得：
\end{enumerate}

\begin{enumerate}
\def\labelenumi{(\roman{enumi})}
\item
  \(s.(t, -1) = (s + t, -1);\)
\item
  \(s.(t,1) = (t - s,1);\)
\item
  \(s.f_{a}(t) = f_{a}(s + t)\) （对 \(a\geqslant 1\) ）
\end{enumerate}

试证 R 连续地作用于 X，命题 4 的四个条件都满足，但 X/R 不是豪斯多夫的。

\begin{enumerate}
\def\labelenumi{\alph{enumi})}
\setcounter{enumi}{1}
\tightlist
\item
  设 S 是 X 上的等价关系，其类由单点 z 组成（若 \(z \notin D' \cup D''\) ），且形如 \((z, -z)\) （若 \(z \in D'\) 或 \(z \in D''\) ）。设 \(X'\) 是商空间 X/S。关系 S 与作用于 X 上的群 R 相容（§ 2，第 4 节）；试证 R 连续地作用于 \(X'\) ，命题 4 的四个条件都满足，且 \(X'/R\) 是豪斯多夫的，但 R 不逆紧地作用于 \(X'\) 。
\end{enumerate}

\begin{enumerate}
\def\labelenumi{\arabic{enumi})}
\setcounter{enumi}{6}
\item
  设 G 是一个豪斯多夫群，连续地作用于一个局部紧空间 X，并设：(i) 对 X 的每个紧子集 K，限制于 \(G \times K\) 的映射 \((s, x) \to s.x\) 是 \(G \times K\) 到 X 中的一个闭映射；(ii) 对每个 \(x \in X\) ，\(s \to s.x\) 是 G 到 X 中的一个逆紧映射。试证 G 逆紧地作用于 X（利用命题 12）。
\item
  设 \(G_{0}\) 是一个非离散的豪斯多夫拓扑群，其中每个紧子集都是有限的（§ 2，习题 23）。设 G 是赋以离散拓扑的群 \(G_{0}\) 。试证 G 经左平移连续但非逆紧地作用于 \(G_{0}\) ，且使得 P(K, L) 对 \(G_{0}\) 的所有紧子集 K、L 都是紧的（定理 1）。
\item
  设 G 是一个拓扑群，连续地作用于两个空间 X 与 X' ，并设 \(f: \mathbf{X} \to \mathbf{X}'\) 是一个连续映射，它与 G 的恒等映射相容（§ 2，第 4 节）。设 \(\vec{f}: \mathbf{X}/\mathbf{G} \to \mathbf{X}'/\mathbf{G}\) 是 \(f\) 在商空间上诱导的连续映射。
\end{enumerate}

\begin{enumerate}
\def\labelenumi{\alph{enumi})}
\item
  试证：若 \(f\) 是开的（相应地，闭的、逆紧的、单射的、满射的），则 \(\overline{f}\) 也是。（为证明若 \(f\) 是逆紧的则 \(\overline{f}\) 是逆紧的，试证 \(\mathbf{X}' / \mathbf{G}\) 的一点在 \(\overline{f}\) 下的逆像是拟紧的。）
\item
  设 G 在 X 与 X' 上都逆紧且自由地作用。试证：若 \(\overline{f}\) 是逆紧的，则 f 也是{[}注意 f 在轨道 G.x 上的限制是 G. \(f(x)\) 上的一个同胚{]}。
\end{enumerate}

¶ 10) 设 G 是一个拓扑群，X、Y 是 G 连续地作用于其上的两个拓扑空间。则 G 按律 (s, (x, y)) → (s.x, s.y) 连续地作用于 X × Y。用 \(X \times {}^{G}Y\) 表示 (X × Y)/G。

\begin{enumerate}
\def\labelenumi{\alph{enumi})}
\item
  试证 \(X \times {}^{G}G\) （G 经左平移作用于自身）典范地同胚于 X。（注意若 U 在 X 中开，则 \(U \times \{e\}\) 的诸点的轨道之并在 \(X \times G\) 中是开的。）
\item
  设 \(Y'\) 是第三个拓扑空间，G 连续地作用于其上，并设 \(f: Y \to Y'\) 是一个连续映射，它与 G 的恒等映射相容。映射 \(\mathrm{I}_{X} \times f\) 诱导出连续映射 \(\bar{f}: X \times {}^{G}Y \to X \times {}^{G}Y'\) 。若 \(f\) 是开的（相应地，逆紧的、单射的、满射的），则 \(\bar{f}\) 也是（习题 9）。给出一个例子，其中 \(f\) 是闭的，\(G\) 逆紧地作用于 \(X\) ，但 \(\bar{f}\) 不是闭的{[}取 \(X = G\) ，\(X' = P\) （由单点组成的空间），并用习题 4{]}。
\item
  试证：若 G 逆紧地作用于 X，则 G 逆紧地作用于 \(X \times Y\) 。
\end{enumerate}

\(d)\) 试证：若 \(\mathbf{Y}\) 是紧的，则典范映射 \(\mathbf{X} \times^{\mathbf{G}}\mathbf{Y} \to \mathbf{X}/\mathbf{G}\) 是逆紧的{[}利用 \(b)\) {]}。

\begin{enumerate}
\def\labelenumi{\alph{enumi})}
\setcounter{enumi}{4}
\tightlist
\item
  设 G 是局部紧的，并设 \(G'\) 是 G 的单点紧化。若 \(\omega\) 是 \(G'\) 中的无穷远点，则 G 按外律连续地作用于 \(G'\) ，该外律由 s.t = st （若 \(t \in G\) ）与 \(s.\omega = \omega\) 给出；于是 X 同胚于 \(X \times {}^{G}G'\) 的一个开子空间（利用 b)）。推断：若 G 逆紧地作用于 X，且此外 X/G 是局部紧的，则 X 是局部紧的{[}利用 c) 证明 \(X \times {}^{G}G'\) 是豪斯多夫的{]}。
\end{enumerate}

\(^{*}f)\) 给出一个局部紧群 G 连续地作用于一个非局部紧空间 X 的例子，使得 X/G 由单点组成（参看 § 2，习题 29）。 *

¶ 11) 设 G 是一个拓扑群，H、K 是 G 的两个闭子群。

\begin{enumerate}
\def\labelenumi{\alph{enumi})}
\item
  试证下列三个条件等价：(i) H × K 按外律 \(((h, k), s) \rightarrow hsk^{-1}\) 逆紧地作用于 G ；(ii) H 逆紧地作用于齐性空间 G/K ；(iii) K 逆紧地作用于齐性空间 G/H。
\item
  若 G 是局部紧的，则 a) 的（等价的）条件也等价于下列条件：对 G 的每对紧子集 A、B，交 HA ∩ BK 是紧的。特别地，若子群 H、K 之一是紧的，则这些条件满足。
\item
  设 G 是局部紧的，试证下列条件等价：(i) 对每个 \(x \in G\) ，映射 \(h \to h.xK\) （H 到 G/K 中的）是逆紧的；(ii) 对每个 \(x \in G\) ，映射 \(k \to k.xH\) （K 到 G/H 中的）是逆紧的；(iii) 对每个 \(x \in G\) 与 G 的每个紧子集 A，\(Hx \cap AK\) 是紧的；(iv) 对每个 \(x \in G\) 与 G 的每个紧子集 A，\(Kx \cap AH\) 是紧的。若子群 H、K 之一是正规的，或者 G 上的左一致结构与右一致结构相同（§3，习题 3），试证这些条件蕴含 a) 的条件。
\end{enumerate}

\begin{enumerate}
\def\labelenumi{\arabic{enumi})}
\setcounter{enumi}{11}
\tightlist
\item
  \begin{enumerate}
  \def\labelenumii{\alph{enumii})}
  \tightlist
  \item
    设 X 是一个紧空间，G 是连续地作用于 X 的一个拓扑群。设每个轨道 A （关于 G 的）相对于其闭包 \(\overline{A}\) 有非空内部。试证至少有一个紧轨道（考虑 X 的在 G 下稳定的紧子集组成的集合中的一个极小元）。
  \end{enumerate}
\end{enumerate}

* b) 给出一个局部紧群连续地作用于一个紧空间的例子，使得任何轨道都不是紧的（§ 2，习题 29）。 *

¶ 13) 设 \(G\) 是一个局部紧群，\(D\) 是 \(G\) 的一个离散子群，使得齐性空间 \(G / D\) 是紧的。试证：对每个 \(d \in D\) ，所有 \(sds^{-1}\) （当 \(s\) 取遍 \(G\) ）组成的集合在 \(G\) 中是闭的。（先借助命题 10 证明：存在一个紧子集 \(K\) （\(G\) 的），使得 \(G = K.D\) ；然后用命题 1 的推论 1。）

\begin{enumerate}
\def\labelenumi{\arabic{enumi})}
\setcounter{enumi}{13}
\tightlist
\item
  \begin{enumerate}
  \def\labelenumii{\alph{enumii})}
  \tightlist
  \item
    设 G 是一个拓扑群，连续地作用于一个拓扑空间 X，并设 \(x_0\) 是 X 的一点，使得 \(s.x_0 = x_0\) 对一切 \(s \in G\) 成立。试证：对 \(x_0\) 的每个邻域 V 与 G 的每个拟紧子集 K，集合 \(\bigcap_{s \in \mathbf{K}} s.V\) 是 \(x_0\) 的一个邻域。
  \end{enumerate}
\end{enumerate}

\begin{enumerate}
\def\labelenumi{\alph{enumi})}
\setcounter{enumi}{1}
\tightlist
\item
  由 a) 推断：若 G 是一个局部紧群，V 是 e 的一个邻域，K 是 G 的一个紧子集，则集合 \(\bigcap_{s\in K}sVs^{-1}\) 是 e 的一个邻域。
\end{enumerate}

¶ 15) 设 G 是一个局部紧群，\(G_{0}\) 是 G 中的单位分量；设商群 \(G/G_{0}\) 是紧的。

\begin{enumerate}
\def\labelenumi{\alph{enumi})}
\item
  试证 G 是 \(\sigma\) -紧的（注意 \(G_{0}\) 是 \(\sigma\) -紧的，并用命题 10）。
\item
  设 \((\mathbf{U}_{n})\) 是 G 中 e 的邻域的一个递减序列。试证存在 G 的一个紧正规子群 K，含于诸 \(U_{n}\) 的交，且使得 G/K 的单位元有一个可数的邻域基本系。{[}存在 e 在 G 中的一个紧对称邻域 \(V_{0}\) ，使得 \(G = V_{0}G_{0}\) 。递归地定义 e 的紧对称邻域的一个序列 \(V_{n}\) ，使得 \(V_{n}^{2} \subset V_{n-1} \cap U_{n}\) 且 \(xV_{n}x^{-1} \subset V_{n-1}\) 对每个 \(x \in V_{0}\) 成立（习题 14 b)；试证诸 \(V_{n}\) 的交 K 回答了问题{]}。
\item
  设 X 是一个局部紧空间，具有一个可数基；设 G 连续地作用于 X，使得 G 中没有 \(s \neq e\) 使 s.x = x 对一切 \(x \in X\) 成立。试证 G 的单位元有一个可数的邻域基本系。{[}设 \((W_n)\) 是 X 的拓扑的一个可数基，由相对紧的集组成；对每对整数 \((m, n)\) （满足 \(\overline{W_m} \subset W_n\) ），设 \(U_{mn}\) 是所有 \(s \in G\) （满足 \(s \cdot \overline{W_m} \subset W_n\) ）组成的集合；试证诸 \(U_{mn}\) 在 G 中开，并应用 b) 的结果。{]}
\end{enumerate}

\begin{enumerate}
\def\labelenumi{\arabic{enumi})}
\setcounter{enumi}{15}
\tightlist
\item
  设 G 是一个局部紧连通群，K 是 G 的一个紧正规子群，N 是 K 的一个闭正规子群，使得 K/N 没有任意小的子群（§ 2，习题 30）。试证在此情形 N 是 G 的一个正规子群。（关于 K/N 的假设蕴含存在 e 在 G 中的一个紧邻域 U ，使得 \(xNx^{-1} \subset N\) 对每个 \(x \in U\) 成立。）
\end{enumerate}

¶ 17) 设 G 是一个没有任意小的子群的局部紧群{[}§ 2，习题 30{]}。

\begin{enumerate}
\def\labelenumi{\alph{enumi})}
\item
  试证 G 的每个点都有一个可数的邻域基本系（习题 15 b)）。
\item
  存在 e 的一个紧邻域 V ，使得对 V 的任意两点 x、y，关系 \(x^{2}=y^{2}\) 蕴含 x=y。（可以设 G 不是交换的。若结论不真，则存在 G 的点的两个序列 \((x_{n})\) 、\((y_{n})\) ，趋于 e，使得 \(x_{n}^{2}=y_{n}^{2}\) 且 \(x_{n}y_{n}^{-1}=a_{n}\neq e\) 。设 U 是 e 的一个紧邻域，它不含 G 的除 \(\{e\}\) 外的任何子群，并设 \(p_{n}\) 是最小的整数 p\textgreater0 使得 \(a_{n}^{p+1}\notin U\) ：试证（必要时过渡到子序列）可以设 \(a=\lim_{n\to\infty}a_{n}^{p_{n}}\) 存在、不等于 e 且属于 U。试证 \(a^{-1}=a\) ，由此得出矛盾）。
\item
  设 U 是 e 的一个紧对称邻域，它不含 G 的除 \(\{e\}\) 外的任何子群，并设 V 是 e 的一个邻域。试证存在一个数 \(c(V) > 0\) ，使得对每对整数 p \textgreater{} 0、q \textgreater{} 0 （满足 \(p \leqslant c(V)q\) ），以及每个元素 \(x \in G\) （满足 \(x, x^{2}, \ldots, x^{q}\) 属于 U），有 \(x^{p} \in V\) 。{[}用反证法，设存在整数的两个序列 \((p_{n}), (q_{n})\) ，使得 \(\lim_{n \to \infty} (p_{n}/q_{n}) = 0\) ，且对每个 n 有一个元素 \(a_{n} \in G\) ，使得 \(a_{n}^{h} \in U\) （对 \(1 \leqslant h \leqslant q_{n}\) ），但 \(a_{n}^{p_{n}} \notin V\) ；还可以设序列 \((a_{n}^{p_{n}})\) 有一个极限 \(a \neq e\) ，它属于 U。试证这时将有 \(a^{m} \in U\) 对每个整数 m \textgreater{} 0 成立，从而得出矛盾{]}。
\end{enumerate}

\begin{enumerate}
\def\labelenumi{\arabic{enumi})}
\setcounter{enumi}{17}
\tightlist
\item
  设 G 是一个局部紧的全不连通拓扑群，其左一致结构与右一致结构相同。试证 e 在 G 中的每个邻域都包含 G 的一个紧的开正规子群（利用第 6 节命题 14 的推论 1，以及 § 3 的习题 3）。
\end{enumerate}

¶ 19) 设 G 是一个局部紧群，\(G_{0}\) 是 G 的单位分量，H 是 G 的一个闭子群。试证：若齐性空间 G/H 是局部连通的，则 \(G_{0}H\) 在 G 中开。设 \(H_{0}\) 是 H 的单位分量；则存在 H 的一个开子群 L ，使得 \(L/H_{0}\) 是紧的（第 6 节，命题 14 的推论 1）；试证一方面 \(LG_{0}\) 在 G 中是闭的（利用命题 1 的推论 1），另一方面 \(LG_{0}\) 在 G 中是开的（考虑 \(G_{0}\) 在 G/L 中的典范像，并用命题 14 的推论 3 以及局部连通空间的商空间是局部连通的这一事实）{[}第一章，§ 11，第 6 节，命题 12){]}。

* 20) 设 \(\varphi\) 是典范同态 \(\mathbf{R} \to \mathbf{T}\) ，\(\theta\) 是一个无理数。用下列律在拓扑空间 \(\mathbf{R}^2 \times \mathbf{T}^2\) 上定义一个群结构

\[
\begin{array}{l} \left(x _ {1}, x _ {2}, t _ {1}, t _ {2}\right) \left(x _ {1} ^ {\prime}, x _ {2} ^ {\prime}, t _ {1} ^ {\prime}, t _ {2} ^ {\prime}\right) \\ = \left(x _ {1} + x _ {1} ^ {\prime}, x _ {2} + x _ {2} ^ {\prime}, t _ {1} + t _ {1} ^ {\prime} + \varphi \left(x _ {2} x _ {1} ^ {\prime}\right), t _ {2} + t _ {2} ^ {\prime} + \varphi \left(\theta x _ {2} x _ {1} ^ {\prime}\right)\right). \end{array}
\]

于是 G 成为一个局部紧拓扑群（甚至是李群）。试证 G 的换位子群不是闭的。 *

\begin{enumerate}
\def\labelenumi{\arabic{enumi})}
\setcounter{enumi}{20}
\tightlist
\item
  设 M 是赋以半群结构的一个紧空间，使得 (i) 映射 \((x,y)\to xy\) （\(M\times M\) 到 M 中的）是连续的；(ii) 对每个 \(a\in M\) ，关系 ax=ay 蕴含 x=y。
\end{enumerate}

\begin{enumerate}
\def\labelenumi{\alph{enumi})}
\item
  试证：若 \(\mathbf{F}\) 是 \(\mathbf{M}\) 的一个闭子集，且 \(x \in \mathbf{M}\) 满足 \(xF \subset F\) ，则 \(xF = F\) 。{[}设 \(y\) 是序列 \((x^n)_{n \geqslant 1}\) 在 \(\mathbf{M}\) 中的一个聚点；试证 \(yF = \bigcap_{n \geqslant 1} x^n F\) ，并推出 \(yxF = yF\) {]}。
\item
  由 a) 推断：若此外对每个 \(a \in M\) ，关系 xa = ya 蕴含 x = y，则 M 是一个紧拓扑群。（先证明存在一个单位元 e；为证 \(x \to x^{-1}\) 是连续的，用反证法，考虑 M 上一个收敛于 e 的超滤子）。
\end{enumerate}

¶ 22) a) 设 G 是一个豪斯多夫拓扑群。试证 G 的每个稳定子集 S （它是紧的，或者是开的且相对紧的）是 G 的一个子群{[}利用上面的习题 21，以及 § 2 的习题 28 a){]}。推断一个紧交换群是一个根群{[}§ 2，习题 28 d){]}。

\begin{enumerate}
\def\labelenumi{\alph{enumi})}
\setcounter{enumi}{1}
\item
  由 a) 推断：紧群 K 的每个局部紧的稳定子集是 K 的一个子群。
\item
  设 G 是一个豪斯多夫拓扑群，S 是 G 的一个稳定的闭子集，具有下列性质：存在 G 的一个子集 K ，它相对于 G 的右一致结构是准紧的，且使得 SK = G。试证 S 是 G 的一个子群。（给定 \(s \in S\) ，递归地定义元素的一个序列 \(x_{i} \in K\) 与元素的一个序列 \(s_{i} \in S\) ，使得 \(s^{-1}x_{i} = s_{i+1}x_{i+1}\) ；注意若 i \textless{} j 则 \(s^{-1}x_{i}x^{-1}_{j} \in S\) ，且对 e 在 G 中的每个邻域 V ，存在一对元素 \(x_{i}, x_{j}\) ，使得 i \textless{} j 且 \(x_{i}x^{-1}_{j} \in V\) ）。
\end{enumerate}

\begin{enumerate}
\def\labelenumi{\arabic{enumi})}
\setcounter{enumi}{22}
\tightlist
\item
  设 \(p\) 是一个素数，\((\mathbf{G}_n)\) 是拓扑群的一个无限序列，每个都等同于离散群 \(\mathbf{Z}/(p^2)\) ，并设 \(\mathbf{H}_n\) 是子群 \((p) / (p^2)\) （\(\mathbf{G}_n\) 的）。设 \(\mathbf{G}\) 是诸 \(\mathbf{G}_n\) 相对于诸 \(\mathbf{H}_n\) 的局部直积（ \(\S 2\) ，习题 26）；则 \(\mathbf{G}\) 是一个局部紧群但不是紧的。试证连续同态 \(u: x \to px\) （\(\mathbf{G}\) 到其自身的）不是 \(\mathbf{G}\) 到 \(u(\mathbf{G})\) 上的严格态射，且 \(u(\mathbf{G})\) 在 \(\mathbf{G}\) 中不闭。
\end{enumerate}
% semantic-fragments: legacy-input-seam
\relax{}
